\section{Related Work}\label{s:related}

Previous studies show conflicting algorithm rankings under various conditions. For example, Mao et al. discovered that reinforcement-learning-based Pensieve outperformed traditional ABR algorithms such as MPC-HM, RobustMPC-HM, and BBA \cite{pensieve-sigcomm}, but Yan et al. observed differently in real-world trials \cite{yan2020-situ}. Hybrid schemes have also outperformed pure strategies. Stick by Huang et al. outperformed existing ABR algorithms by 9.4\%–44.3\% in average QoE under varying network conditions \cite{stick-infocom20}, while DeepMPC similarly outperformed Pensieve by approximately 4–5.9\% QoE \cite{deepmpc-icip20}. Narayanan et al.’s 5G measurements revealed that under mmWave conditions, RobustMPC was the least stalling ABR algorithm (less than 5\%), while algorithms such as Pensieve and FastMPC experienced increases in stall times of approximately 82\% and 259\% in 5G \cite{5g-variegated}. Pensieve experienced more stalls on 5G while being the best in 4G traces \cite{5g-variegated}, reflecting the similarity to our own results. These findings also align with CausalSim’s conclusion that simulator bias can change rankings, reducing error by approximately 50–60\% and producing different insights into ABR algorithms than biased simulators \cite{causalsim-nsdi23}. These comparisons highlight that the ABR algorithm ranking reported is highly context-dependent. We therefore replay our six ABR algorithms under both Mahimahi and CellReplay to observe to what extent this is true for the choice of emulator. We did not observe differences in the ranking, but we observed some improvements when CellReplay is used. In general, the performance of ABR algorithms varies significantly with environment, trace complexity, and evaluation setup. Our results extend this understanding by showing that emulator choice and trace variability can significantly affect the performance of ABR algorithms. 
Hybrid or rule-based methods remain strong and balanced across setups, whereas learned policies such as Pensieve must be carefully tuned for each environment.

Comparing our CellReplay‐based rankings with previous work reveals some deviations and similarities. For example, BOLA in its original evaluation achieved 84–95\% of its ideal performance~\cite{bola-bola}, but in our Mahimahi-based experiments, BOLA became the weakest algorithm in 4G and became the least stalling algorithm in 5G settings. This suggests that network conditions can change rankings. Similarly, SENSEI’s user-sensitivity model shows that weighting quality by viewer importance can change results. Their evaluation reports a roughly 15\% QoE gain by adapting to content sensitivity~\cite{sensei-nsdi21}, implying that rankings may change under personalized QoE metrics. Finally, new codec-aware systems such as Swift and LiFteR shift adaptation to the coding layer. Swift’s layered neural codec~\cite{swift-nsdi22} and LiFteR's loose frame reference (parallel learned codec decoding) achieve much higher frame rates and quality~\cite{lifter-nsdi24}. These innovations suggest future ABR algorithms may choose which layers or frames to fetch, potentially altering relative performance. In summary, while our general trends agree with previous studies, using CellReplay and accounting for user preferences or new codec designs can reveal different rankings and merit further investigation.

%%% Local Variables:
%%% mode: latex
%%% TeX-master: "../thesis"
%%% End:
